

\subsubsection{偶数分频（占空比50\%）}


在输入时钟的驱动下，一次输出翻转只对应
$0\rightarrow 1$ 或 $1\rightarrow 0$，
两次翻转才构成一个完整周期。

因此，只需要通过计数器数拍，并在计满半个输出周期时
翻转输出，即可实现偶数分频。

例如，每经过 $5$ 个输入时钟周期翻转一次输出，
则一个完整的输出周期包含 $10$ 个输入时钟周期，
因此实现了 $10$ 分频。

设分频系数为偶数 $N$，则有
\[
T_{\mathrm{out}} = N T_{\mathrm{in}},
\qquad
f_{\mathrm{out}} = \frac{f_{\mathrm{in}}}{N}.
\]

输出每隔 $\frac{N}{2}$ 个输入时钟周期翻转一次，即
\[
T_{\mathrm{toggle}} = \frac{N}{2}T_{\mathrm{in}}.
\]

对应到代码上（以 6 分频为例），则需要我们
复位解除后，计数器从 $0$ 开始计数。
输入时钟上升沿到来时，执行如下动作：

\begin{table}[H]
  \centering
  \caption{6 分频的计数与翻转过程}
  \begin{tabularx}{\textwidth}{ccX}
    \toprule
    输入上升沿 & 沿到来前的计数值 & 该沿执行的动作 \\
    \midrule
    第一个 & 0 & 计数值更新为 1。 \\
    第二个 & 1 & 计数值更新为 2。 \\
    第三个 & 2 & 已计满 3 拍，输出翻转，同时计数器清零。 \\
    \bottomrule
  \end{tabularx}
\end{table}

因此，当计数器从 $0$ 开始计数时，可使用以下条件判断翻转时机：
\[
\boxed{\texttt{count} = \frac{N}{2}-1}
\]

对于 $6$ 分频，判断阈值为
\[
\texttt{count} = \frac{6}{2}-1 = 2.
\]

\subsubsection{奇数分频（ 50\% 占空比）}

如果只使用输入时钟的上升沿，或者只使用输入时钟的下降沿，
则只能按照整数个输入时钟周期安排输出的高、低电平持续时间。

当分频系数为奇数时，若要求占空比为 $50\%$，
则输出高电平与低电平的持续时间均为
\[
T_{\mathrm{high}} = T_{\mathrm{low}}
= \frac{N}{2}T_{\mathrm{in}}.
\]

由于 $N$ 为奇数，$\frac{N}{2}$ 不是整数，
因此需要能够表示半个输入时钟周期的方法。

例如，$5$ 分频且占空比为 $50\%$ 时，需要满足
\[
T_{\mathrm{high}} = T_{\mathrm{low}}
= 2.5T_{\mathrm{in}}.
\]

一种方法是同时利用输入时钟的上升沿和下降沿。
当输入时钟占空比为 $50\%$ 时，相邻上升沿与下降沿之间
恰好相隔半个输入时钟周期。

可以将最终输出拆分为两个信号 $A$、$B$ 的或运算。
其中，$B$ 相对于 $A$ 延迟半个输入时钟周期。

以 5 分频为例的波形组合，以输入时钟周期为时间单位，其高电平区间如下：

\begin{table}[H]
  \centering
  \caption{5 分频的高电平区间组合}
  \begin{tabular}{cc}
    \toprule
    信号 & 高电平区间 \\
    \midrule
    $A$ & $0 \sim 2$ 拍 \\
    $B$ & $0.5 \sim 2.5$ 拍 \\
    $A \lor B$ & $0 \sim 2.5$ 拍 \\
    \bottomrule
  \end{tabular}
\end{table}

两个信号的高电平区间在中间部分重叠，
只在起始与结束处各相差半拍。
通过或运算，合成信号的高电平持续时间为
\[
T_{\mathrm{high}} = \frac{N}{2}T_{\mathrm{in}}.
\]

合成输出可以表示为
\[
\boxed{\texttt{clk\_out} = A \lor B}
\]

在代码实现上：
分别利用上升沿和下降沿更新信号 $A$、$B$，
并提出在以下计数条件下翻转：
\[
\texttt{counter} = \frac{N-1}{2}-1.
\]


\subsubsection{时钟IP核}
时钟IP虽然可以提供准确的时钟，但是需要一段时间来建立稳定的时钟关系。
当状态信号$ locked = 1 $时，时钟才正式准备好。


\subsubsection{testbench仿真平台}

仿真时长需要根据待验证信号的周期合理设置。

如果仿真时间不足以覆盖待验证信号的变化过程，
就可能无法观察到预期现象。因此，应保证仿真时间能够覆盖
复位过程以及若干个完整的输出时钟周期。